%!TEX root main.tex
\section{Linearly-finite series and recognisable series}
\label{sec:recognisable}
% lax begin Lax619925.Recognisable

This brief section does not contain any new results,
however it will be instructive in comparison with the stronger results
developed in~\cref{sec:Hadamard,sec:shuffle,sec:infiltration}.
%
%Recall that a \emph{vector space} over the rationals $\Q$ is a tuple $\V = \tuple{V; %\zero, \cdot, +}$
%consisting of the zero vector $\zero$, a scalar product $\cdot : \Q \times V \to V$a
%and compatible addition operation $+ : V \times V \to V$.
%
% A \emph{linear transition system} consists of a set $V$ simultaneously carrying the structure of a weighted transition system and a vector space.
% %
% Formally, a \emph{linear transition system} (\LTS) is a tuple
% $\A = \tuple{\Sigma, V; F, \delta; \zero, \cdot, +}$
% where $\T = \tuple {\Sigma, V; F, \delta}$ is a \WTS\-
% and $\V = \tuple{V; \zero, \cdot, +}$ is a vector space (over the rationals $\Q$).
% %
% Moreover, we require that the output mapping $F : V \to \Q$
% and the transition functions $\delta_a : V \to V$ (for all $a \in \Sigma$) are linear.
% %
% For instance, series with left derivatives
% $\Tleft = \tuple{\Sigma, \series \Q \Sigma; \e, \deriveleft {}; \zero, \cdot, +}$
% form a linear transition system.

\paragraph{Linearly-finite series.}

We recall a semantic mechanism to define classes of series
that will be useful in the rest of the paper in comparison to more general classes.
%
Given a set of series $F \subseteq \series \Q \Sigma$,
let their \emph{linear span}, denoted by $\linearspan \Q F$,
be the set of finite linear combinations of series from $F$,
with coefficients from $\Q$.
%
In other words, elements of this set are of the form $c_1 \cdot f_1 + \cdots + c_n \cdot f_n$
for some $f_1, \dots, f_n \in F$ and $c_1, \dots, c_n \in \Q$.
%
It is easy to see that $\linearspan \Q F$ is a vector space over $\Q$ containing $F$,
and in fact it is the smallest such vector space.
%
Sets of series of the form $\linearspan \Q {f_1, \dots, f_k}$ are called \emph{finitely generated},
with \emph{generators} $f_1, \dots, f_k$.
%
We say that a series $f \in \series \Q \Sigma$ is \emph{linearly finite}
if it belongs to a finitely generated $\Q$-vector space of series
closed under left derivatives $\deriveleft a$, $a \in \Sigma$.
%
In other words, $f$ is linearly finite when there are generators $f_1, \dots, f_k \in \series \Q \Sigma$ \st:
%
\begin{enumerate}[label=\textbf{(R.\arabic*)}]
    \item \label{linearly_finite:1} $f \in \linearspan \Q {f_1, \dots, f_k}$ and
    \item \label{linearly_finite:2} $\deriveleft a f_i \in \linearspan \Q {f_1, \dots, f_k}$
    for every $a \in \Sigma$ and $1 \leq i \leq k$,
\end{enumerate}

\begin{remark}
    We would get the same class if we additionally required that the generators
    are obtained from $f$ itself by a finite number of left derivatives.
    %
    In other words, we could take the generators to be of the form $f_i = \deriveleft {w_i} f$ for some words $w_i \in \Sigma^*$.
    %
    Since this will not be the case for the more general classes of series that we will consider in the rest of the paper,
    we have chosen the definition above.
\end{remark}

\begin{lemma}
    \label{lem:recognisable closure properties}
    The class of linearly finite series is effectively closed under
    scalar product, addition, left derivatives and anti-derivatives.
\end{lemma}
%
\begin{proof}
    Closure under scalar product follows from the definition of linear span.
    %
    Regarding closure under addition, let $f, g$ be linearly finite series with generators $F$, resp., $G$.
    Then $f + g$ is linearly finite with generators $F \cup G$.
    %
    Closure under left derivatives holds by design:
    If $f$ is linearly finite with generators $G$,  %= \set {f_1, \dots, f_k}$,
    then for every $a \in \Sigma$ the series $\deriveleft a f$ is also linearly finite with the same set of generators.
    %
    % Indeed, by \cref{linearly_finite:1} $f$ is of the form $\sum_{i = 1}^k c_i \cdot f_i$ for some $c_i \in \Q$,
    % and thus by linearity of left derivative we have $\deriveleft a f = \sum_{i = 1}^k c_i \cdot \deriveleft a f_i$.
    % %
    % But now by \cref{linearly_finite:2} we can replace each $\deriveleft a f_i$ by a linear combination of the generators,
    % showing that $\deriveleft a f$ is in the linear span of $G$.
    %
    Finally, for closure under anti-derivatives
    consider linearly finite series $f_1, \dots, f_d$ with generators, resp., $G_1, \dots, G_d$
    \st~$\deriveleft {a_j} f = f_j$ for every $1 \leq j \leq d$.
    %
    Then $f$ is linearly finite with generators $\set {f, f_1, \dots, f_d} \cup G_1 \cup \cdots \cup G_d$.
\end{proof}

The class of linearly finite series is also closed under right derivatives,
however we will obtain this from a stronger closure property, namely closure under reversal.
%
To show the latter, we will use a stronger representation for linearly finite series,
which we recall next.

\paragraph{Linear representations.}

A \emph{linear representation} is a tuple $\tuple{\Sigma, k, x, y, M}$,
where $\Sigma = \set{a_1, \dots, a_d}$ is a finite alphabet,
$k \in \N$ is the \emph{dimension} of the representation,
$x \in \Q^{1 \times k}$ is a row vector of \emph{initial weights},
$y \in \Q^{k \times 1}$ is a column vector of \emph{final weights},
and $M : \Sigma \to \Q^{k \times k}$ is the \emph{transition function},
assigning to each input letter $a \in \Sigma$ a transition matrix $M(a) \in \Q^{k \times k}$.
%
The transition function $M$ is extended homomorphically to all input words $M : \Sigma^* \to \Q^{k \times k}$
by $M(\e) := I_k$ (the $k \times k$ identity matrix)
and $M(a \cdot w) := M(a) \cdot M(w)$ for every $a \in \Sigma$ and $w \in \Sigma^*$.
%
A linear representation \emph{recognises} the series $f \in \series \Q \Sigma$
mapping an input word $w \in \Sigma^*$ to $\coefficient w f := x \cdot M(w) \cdot y$.
%
A series is \emph{recognisable} if it is recognised by some linear representation.
%
The following coincidence result is classical (\cf~\cite[Proposition~5.1]{BerstelReutenauer:CUP:2010}).
%
\begin{lemma}
    A series is recognisable if, and only if, it is linearly finite.
\end{lemma}
%
% \begin{proof}
%     For the ``only if'' direction, assume that $f$ is recognisable.
%     Let $\tuple{\Sigma, k, x, y, M}$ be a linear representation recognising $f$.
%     %
%     Consider the series $g := M \cdot y \in \series {\Q^{k \times 1}} \Sigma$.
%     %
%     In other words, $g_i(w)$ is the $i$-th component of the vector $M(w) \cdot y$.
%     %
%     Then $f$ is linearly finite with generators $g_1, \dots, g_k$.
%     Indeed, $f = u \cdot M \cdot y = \sum_{i = 1}^k u_i \cdot g_i$, showing~\cref{linearly_finite:1}.
%     Moreover, $\deriveleft a g \in \series {\Q^{k \times 1}} \Sigma$ is the series that maps an input word $w$ to
%     %
%     \begin{align*}
%         g(a \cdot w)
%         = M(a) \cdot M(w) \cdot y
%         = M(a) \cdot g(w),
%     \end{align*}
%     %
%     and thus all components of $\deriveleft a g = M(a) \cdot g$ are in the linear span of the generators.
%
%     For the other direction, assume that $f$ is linearly finite
%     with generators $g = \tuple{g_1, \dots, g_k} \in \series \Q \Sigma^{k \times 1}$.
%     %
%     Since $f$ is in the linear span of the generators by~\cref{linearly_finite:1},
%     there are coefficients $x \in \Q^{1 \times k}$ \st~$f = x \cdot g$.
%     %
%     For every input symbol $a \in \Sigma$,
%     by~\cref{linearly_finite:2} each component of $\deriveleft a g$ is in the linear span of the generators,
%     and thus there is a matrix $M(a) \in \Q^{k \times k}$ \st~$\deriveleft a g = M(a) \cdot g$.
%     %
%     Finally, let $y = g(\e) \in \Q^{k \times 1}$ be the vector of constant terms of the generators.
%     %
%     Then $f$ is recognised by the linear representation $\tuple{\Sigma, k, x, y, M}$.
%     %
%     \begin{claim*}
%         We have $g = M \cdot y$.
%     \end{claim*}
%     \begin{claimproof}
%         We proceed by coinduction.
%         In the base case, we have $g(\e) = y = M(\e) \cdot y$.
%         %
%         In the coinductive case we have 
%         $\deriveleft a g = M(a) \cdot g = M(a) \cdot M \cdot y = (\deriveleft a M) \cdot y = \deriveleft a (M \cdot y)$.
%     \end{claimproof}
%     %
%     Thanks to the claim we have $f = x \cdot g = x \cdot M \cdot y$,
%     as required.
% \end{proof}
%
\noindent
Thanks to the lemma above,
recognisable series are closed under the operations from~\cref{lem:recognisable closure properties}.
%
The advantage of linear representations is that they make it possible to show stronger closure properties,
as we now demonstrate.
%
The \emph{reversal} of series $f \in \series \Q \Sigma$ is the series $\reversal f \in \series \Q \Sigma$
which reverses the order of the letters in the input.
%
That is, for every input word $w \in \Sigma^*$ we have $\coefficient w {(\reversal f)} := \coefficient {\reversal w} f$,
where $\reversal w$ is the word obtained from $w$ by reversing the order of its letters.
%
For instance, $\reversal {(a b c)} = c b a$.
%
Note that reversal is an endomorphism on the vector space of series.
%
\begin{lemma}
    If $f \in \series \Q \Sigma$ is recognisable, then $\reversal f$ is also recognisable.
\end{lemma}
%
\begin{proof}
    Let $f$ be recognised by a linear representation $\tuple{\Sigma, k, x, y, M}$.
    Then its reversal $\reversal f$ is recognised by the linear representation $\tuple{\Sigma, k, \transpose y, \transpose x, \transpose M}$,
    where $\transpose M$ is the transpose of $M$, and similarly for $\transpose y, \transpose x$.
    %
    Indeed, for every input word $w \in \Sigma^*$ we have
    %
    $
        \coefficient w {(\reversal f)}
        = \coefficient {\reversal w} f
        = x \cdot M(\reversal w) \cdot y
        = x \cdot \transpose {(\transpose M(w))} \cdot y
        = \transpose y \cdot \transpose M(w) \cdot \transpose x
    $.
\end{proof}
%
\noindent
Notice that it is not so clear how to prove closure under reversal
solely based on the definition of linearly finite series.
%
Since recognisable series are closed under left derivatives and reversal,
by double reversal $\deriveright a f = \reversal {(\deriveleft a \reversal f)}$
we obtain closure under right derivatives as well.
%
\begin{corollary}
    \label{cor:recognisable derive right}
    If $f$ is recognisable then $\deriveright a f$ is also recognisable.
\end{corollary}
% \begin{proof}
    % Let $f_1, \dots, f_k$ be generators for $f$.
    % We adjoin new generators $\deriveright a f_i$
    % for every $a \in \Sigma$ and $1 \leq i \leq k$.
    % %
    % Since $f$ is in the linear span of the $f_i$'s,
    % we have that $\deriveright a f$ is in the linear span of the $\deriveright a f_i$'s,
    % thus the first condition is satisfied.
    % %
    % Now consider $\deriveleft a \deriveright b f_i$.
    % Since left and right derivatives commute by~\cref{lem:derive left right commutativity},
    % we have $\deriveleft a \deriveright b f_i = \deriveright b \deriveleft a f_i$.
    % %
    % But $\deriveleft a f_i$ is in the linear span of the $f_i$'s.
    % Since right derivation is linear,
    % $\deriveright b \deriveleft a f_i$ is in the linear span of the $\deriveright b f_i$'s.
    % \qedhere
% \end{proof}

\begin{theorem}
    The class of recognisable (= linearly finite) series is an effective prevariety.
    In particular, the commutativity problem is decidable.
\end{theorem}

% \paragraph{Minimisation of linear representations.}
% %
% If a rational series is commutative,
% then in a minimal linear representation the transition matrices commute pairwise.

\paragraph{Towards more general classes.}

In the next sections~\cref{sec:Hadamard,sec:shuffle,sec:infiltration}
we will consider classes of series which generalise the recognisable ones.
%
In each case, the more general class will be obtained by considering a notion of product of series,
thus moving from a vector space structure to a more general algebra structure.
%
The basic closure properties as in~\cref{lem:recognisable closure properties} will follow from first principles.
%
Decidability of the zeroness problem will be obtained as an application of Hilbert finite basis theorem
for finitely generated polynomial algebras with finitely many variables.
%
The price to pay is that such classes of series will not be closed under reversal in general.
%
Consequently, in order to show that they are prevarieties
we will need to show closure under right derivatives by other means.
% lax end